\scanpage{11}
\uline{Dualité} \deleted{(\illegible{1})}

La condition $\rho$ inv est néc. et suff pour que l'accouplement naturel entre $V$ et $V'$ soit une dualité -- elle implique donc que $C$ et $C'$ sont en dualité l'un avec l'autre.

Si de plus $\beta$ inv. i.e. \deleted{\unsure{applt}} $f\mid V$ non dég. i.e. $\psi_V:V\simeq V'$, on trouve un iso entre \deleted{$V$ et} $C$ et $C'$ qui commute à $\unsure{\mathscr G|G|\mathfrak G}$ et transforme donc $Q_C$ en $Q_{C'}$. Donc dans ce cas
\[
(C,Q_C)\simeq(C',Q_{C'})
\]

Le passage d'un \deleted{\unsure{hy|hyp}}polyèdre régulier \unsure{purement|projectif} épinglé vers son dual, sous ces condition, en termes de $C^\infty_{pr\,orth}$, s'exprime donc par la conservation de $(C,Q_C)$, et en \deleted{\unsure{échangeant} \unsure{à}} \inserted{remplaçant} $\tau_0^C\ldots\tau_n^C$ par $\tau_n^C\ldots\tau_0^C$ (i.e. $\varphi_g^C$ par $\varphi_g^C\circ\theta$, où $\theta:\mathscr G\longrightarrow\mathscr G$ échange $\underline\tau_i$ et $\underline\tau_{n-i}$) -- ici il est inutile de préciser une section $a_0$ \inserted{resp $a_n$} de $C$, \deleted{\unsure{car}} puisque $\beta_0,\beta_n$ inversibles.

% archive sheet number: 10

